\documentclass[11pt]{article}
\usepackage[margin=1in]{geometry}
\usepackage{booktabs}
\title{Multilayer Perceptron vs Linear Baseline on Handwritten Digits}
\author{Devis W. Saputra}
\begin{document}
\maketitle

\begin{abstract}
This experiment compares a two-hidden-layer MLP with standardized logistic regression on the scikit-learn handwritten-digits benchmark to test whether nonlinear complexity improves held-out classification.
\end{abstract}

\section{Method}
Both models use a stratified 75/25 split with seed 42 and standardized inputs. The MLP contains hidden layers of 128 and 64 units with early stopping.

\section{Results}
\begin{tabular}{lrr}
\toprule
Model & Accuracy & Macro-F1 \\
\midrule
Logistic regression & 0.9778 & 0.9776 \\
MLP & 0.9578 & 0.9575 \\
\bottomrule
\end{tabular}

The MLP stops after 29 optimization iterations.

\section{Interpretation}
The linear baseline outperforms the MLP on this small 8-by-8 benchmark. The negative result is retained because model complexity should be justified empirically rather than assumed to be superior.

\section{Limitations}
Repeated seeds, nested tuning, uncertainty intervals, and larger image datasets would support stronger conclusions.

\section{Calculation audit and reproduction scope}
This experiment compares a multilayer perceptron with multinomial logistic regression on the same standardized handwritten-digit split. The recorded linear baseline reaches 0.9778 accuracy, compared with 0.9578 for the MLP, with macro-F1 showing the same ordering. Keeping this negative result makes the engineering question clear: additional nonlinear capacity must justify its complexity under an explicit evaluation protocol.

Macro-F1 gives each class equal weight. A single seeded holdout does not establish a universal ranking of architectures; training and preprocessing choices remain part of the comparison.

See the repository calculation guide for exact source paths. The full experiment was not independently rerun during this documentation review.
\end{document}
